\centerline{\bf \Huge Вписанные уголки}

\begin{flushright}
        \scriptsize{ну это хоть должны уметь решать...}
\end{flushright}

\problem{1}
Окружности $\omega_1$ и $\omega_2$ пересекаются в точках $K$ и $L$. Прямая $\ell$ касается окружностей $\omega_1$ и $\omega_2$ в точках $X$ и $Y$ соответственно (точка $K$ лежит внутри треугольника $XYL$). Прямая $XK$ вторично пересекает $\omega_2$ в точке $Z$. Докажите, что $LY$ --- биссектриса угла $\angle XLZ$.

\problem{2}
Окружность $\omega$ с центром $O$ и окружность $\omega_1$ пересекаются в точках $A$ и $B$. На дуге окружности $\omega$, лежащей внутри $\omega_1$, взята точка $C$. Вторые точки пересечения прямых $AC$ и $BC$ с $\omega_1$ обозначим через $E$ и $D$ соответственно. Докажите, что прямые $ED$ и $OC$ перпендикулярны.

\problem{3}
Пусть $AD$ --- высота прямоугольного треугольника $ABC$ ($\angle BAC=90^\circ$). Точки $I_1$ и $I_2$ --- центры вписанных окружностей треугольников $ADB$ и $ADC$ соответственно. Окружность с центром $A$ и радиусом $AD$ повторно пересекает катеты треугольника в точках $K$ и $L$. Докажите, что точки $I_1$, $I_2$, $K$ и $L$ лежат на одной прямой.

\problem{4}
Дана равнобокая трапеция $ABCD$ с основаниями $BC$ и $AD$. Окружность $\omega$ проходит через вершины $B$ и $C$ и вторично пересекает сторону $AB$ и диагональ $BD$ в точках $X$ и $Y$ соответственно. Касательная, проведенная к окружности $\omega$ в точке $C$, пересекает луч $AD$ в точке $Z$. Докажите, что точки $X$, $Y$ и $Z$ лежат на одной прямой.

\problem{5}
Отрезки $BB_1$ и $CC_1$ --- высоты остроугольного треугольника $ABC$. Прямая, проходящая через центры вписанных окружностей треугольников $BCC_1$ и $CBB_1$, пересекает стороны $AB$ и $AC$ в точках $X$ и $Y$ соответственно. Докажите, что $AX=AY$.

\problem{6}
Окружность $\omega$ проходит через три середины сторон треугольника $ABC$. Стороны треугольника отсекают от $\omega$ три маленькие дуги. Докажите, что одна из них равна сумме двух других.

\problem{7}
В окружность вписан выпуклый пятиугольник $ABCDE$. Известно, что $AE=DE$. Диагонали $AC$ и $BD$ пересекаются в точке $P$. На продолжении стороны $AB$ за точку $A$ отмечена такая точка $Q$, что $AQ=DP$. На продолжении стороны $DC$ за точку $D$ отмечена такая точка $R$, что $DR=AP$. Докажите, что $PE\perp QR$.

\problem{8}
Точка $I$ --- центр вписанной в треугольник $ABC$ окружности. Внутри треугольника $ABC$ расположена окружность $\omega$, которая касается сторон $AB$ и $AC$ в точках $X$ и $Y$. Пусть $Z$ --- одна из двух точек пересечения $\omega$ с описанной окружностью треугольника $BIC$. Докажите, что описанные окружности треугольников $BXZ$ и $CYZ$ касаются.

\problem{9}
Внутри выпуклого четырёхугольника $ABCD$ нашлась такая точка $P$, что выполняются равенства

$$\angle PAD: \angle PBA:\angle DPA=1:2:3=\angle CBP:\angle BAP:\angle BPC.$$

Докажите, что следующие три прямые пересекаются в одной точке: внутренние биссектрисы углов $\angle ADP$ и $\angle PCB$ и серединный перпендикуляр к отрезку $AB$.